% !TeX program = pdflatex
% Universite de N'Zerekore - Licence 1 L1 Chimie
% Algèbre - 01-MTH 1102 - Chapitre 1 : Ensembles, applications et raisonnement
% Serie d'entrainement de 30 exercices (sans corriges)
\documentclass[11pt,a4paper,oneside]{article}
\usepackage[T1]{fontenc}
\usepackage[utf8]{inputenc}
\usepackage[main=french,provide=*]{babel}
\usepackage{lmodern,microtype}
\usepackage{amsmath,amssymb,mathtools}
\usepackage{geometry}
\geometry{left=2.3cm,right=2.3cm,top=2.55cm,bottom=2.35cm,headheight=24pt}
\usepackage{array,booktabs,tabularx}
\usepackage{enumitem}
\usepackage{xcolor}
\usepackage{tikz}
\usetikzlibrary{calc}
\usepackage[most]{tcolorbox}
\usepackage{fancyhdr,lastpage}
\usepackage{titlesec}
\usepackage{needspace}
\usepackage{hyperref}
\definecolor{Marine}{HTML}{10283F}
\definecolor{Bleu}{HTML}{176D8E}
\definecolor{Sarcelle}{HTML}{118F86}
\definecolor{Or}{HTML}{EDAA47}
\definecolor{GrisTexte}{HTML}{54667A}
\definecolor{BleuPale}{HTML}{EFF5F8}
\definecolor{OrangePale}{HTML}{FFF6E7}
\hypersetup{colorlinks=true,linkcolor=Bleu,urlcolor=Sarcelle,
pdftitle={Algèbre L1 Tronc commun - Chapitre 1 - 30 exercices d'entrainement},
pdfauthor={M. Bakary TRAORE - Universite de N'Zerekore}}
\setlength{\parindent}{0pt}
\setlength{\parskip}{4pt plus1pt minus1pt}
\setlist[enumerate]{leftmargin=1.6em,itemsep=2pt,topsep=3pt}
\setlist[itemize]{leftmargin=1.4em,itemsep=2pt,topsep=3pt}
\renewcommand{\arraystretch}{1.15}
\emergencystretch=2em
\titleformat{\section}{\Large\bfseries\color{Marine}}{\thesection}{.65em}{}[\vspace{1mm}{\color{Sarcelle}\titlerule[.75pt]}]
\titleformat{\subsection}{\normalsize\bfseries\color{Bleu}}{}{0pt}{}
\titlespacing*{\section}{0pt}{18pt}{9pt}
\titlespacing*{\subsection}{0pt}{10pt}{4pt}
\pagestyle{fancy}
\fancyhf{}
\fancyhead[L]{\small\color{Marine}\textbf{UNIVERSITÉ DE N'ZÉRÉKORÉ}}
\fancyhead[R]{\small\color{GrisTexte}Algèbre \textbullet\ 01-MTH 1102}
\fancyfoot[L]{\scriptsize\color{GrisTexte}Chapitre 1 \textbullet\ L1 Chimie}
\fancyfoot[R]{\scriptsize\color{GrisTexte}\thepage\,/\,\pageref{LastPage}}
\renewcommand{\headrulewidth}{.5pt}
\renewcommand{\footrulewidth}{.3pt}
\renewcommand{\headrule}{\hbox to\headwidth{\color{Sarcelle}\leaders\hrule height \headrulewidth\hfill}}
\renewcommand{\footrule}{\hbox to\headwidth{\color{Bleu!30}\leaders\hrule height \footrulewidth\hfill}}
\newtcolorbox{objectifbox}[1][]{enhanced,breakable,colback=BleuPale,colframe=Bleu,
boxrule=.6pt,arc=1mm,#1}
\newtcolorbox{conseilbox}[1][]{enhanced,breakable,colback=OrangePale,colframe=Or!75!black,
boxrule=.6pt,arc=1mm,#1}
\newcommand{\exo}[3]{%
  \Needspace{6\baselineskip}%
  \subsection*{Exercice #1 --- #3 \hfill
  {\normalfont\footnotesize\color{GrisTexte}[#2]}}}
\newcommand{\R}{\mathbb{R}}
\newcommand{\D}{\mathcal{D}}
\begin{document}
\hypersetup{pageanchor=false}
\thispagestyle{empty}
\begin{tikzpicture}[remember picture,overlay]
\fill[Bleu] (current page.south west) rectangle (current page.north east);
\draw[white,opacity=.09,line width=1pt] ([xshift=2.6cm,yshift=5cm]current page.south west) circle (4.1cm);
\draw[white,opacity=.09,line width=1pt] ([xshift=16.2cm,yshift=23.5cm]current page.south west) circle (5.4cm);
\draw[Or,opacity=.8,line width=2pt] ([xshift=3cm,yshift=7.2cm]current page.south west) -- ([xshift=18cm,yshift=7.2cm]current page.south west);
\end{tikzpicture}
\color{white}\vspace*{2.85cm}
\begin{center}
{\fontsize{19}{23}\selectfont\bfseries UNIVERSITÉ DE N'ZÉRÉKORÉ\par}
\vspace{.5cm}{\large LICENCE 1 EN CHIMIE\par}
\vspace{2.1cm}
{\color{Or}\fontsize{20}{25}\selectfont\bfseries FEUILLE D'ENTRAÎNEMENT\par}
\vspace{.75cm}{\fontsize{34}{40}\selectfont\bfseries ALGÈBRE\par}
\vspace{.45cm}{\fontsize{25}{31}\selectfont\bfseries CHAPITRE 1\par}
\vspace{.5cm}{\fontsize{20}{25}\selectfont Ensembles, applications et raisonnement\par}
\vspace{1cm}{\color{Or}\fontsize{24}{29}\selectfont\bfseries 30 EXERCICES PROGRESSIFS\par}
\vspace{.5cm}{\large Ensembles \textbullet\ Applications \textbullet\ Logique \textbullet\ Chimie\par}
\vspace{1.5cm}{\normalsize Code : 01-MTH 1102 \quad | \quad Exercices sans corrigé\par}
\vspace{.5cm}{\normalsize Enseignant : \textbf{M. Bakary TRAORE}\par}
\vspace{.35cm}{\normalsize Année universitaire 2026--2027\par}
\end{center}
\clearpage\color{black}\pagenumbering{arabic}
\hypersetup{pageanchor=true}\setcounter{page}{1}

\section*{Présentation de la série}
\addcontentsline{toc}{section}{Présentation de la série}
\begin{objectifbox}
\textbf{Public :} Licence 1 Chimie --- ECUE Algèbre, UE Mathématiques 1.\par
\textbf{Chapitre 1 :} Ensembles, applications et raisonnement (5 h : 3 h CM + 2 h TD).\par
\textbf{Finalité :} Consolider les notions d'ensemble, de produit cartésien, d'application et de preuve, puis les utiliser pour décrire des situations chimiques simples.\par
\textbf{Consignes :} distinguer élément et sous-ensemble, préciser l'ensemble de départ et d'arrivée des applications, rédiger des preuves courtes et vérifier les propriétés demandées.
\end{objectifbox}
\begin{center}\begin{tabularx}{\textwidth}{@{}cXc@{}}\toprule
\textbf{Série} & \textbf{Thème} & \textbf{Exercices}\\\midrule
1 & Ensembles, appartenance et inclusion & 1--5\\
2 & Opérations sur les ensembles et lois de De Morgan & 6--10\\
3 & Produit cartésien et contraintes de formulation & 11--15\\
4 & Applications, injection, surjection et réciproque & 16--20\\
5 & Logique, quantificateurs et preuves & 21--25\\
6 & Applications aux données et aux contraintes chimiques & 26--30\\
\bottomrule\end{tabularx}\end{center}
\begin{conseilbox}\textbf{Entraînement complémentaire :} travailler les 30 exercices sans consulter les solutions. Cette feuille complète le polycopié de chapitre 1 et ne remplace pas le TD officiel.\end{conseilbox}
\setcounter{tocdepth}{1}\tableofcontents
\section{Ensembles, appartenance et inclusion}

\exo{1}{Facile}{Éléments et parties}
Soit $A=\{\mathrm C,\mathrm H,\mathrm O\}$, où les lettres désignent des éléments chimiques.
Dire si $\mathrm H\in A$, $\{\mathrm H\}\in A$, $\{\mathrm H\}\subseteq A$ et $\mathrm N\in A$ sont vraies ou fausses ; justifier les distinctions.

\exo{2}{Facile}{Égalité d'ensembles}
Pour $A=\{1,2,3\}$ et $B=\{3,1,2\}$, montrer $A=B$ par double inclusion. Expliquer pourquoi l'ordre d'énumération n'intervient pas.

\exo{3}{Facile}{Définition par compréhension}
Écrire en extension l'ensemble $E=\{n\in\mathbb Z:-2\le n\le3\}$.
Écrire ensuite en compréhension l'ensemble $F=\{0,2,4,6,8\}$ des entiers pairs considérés.

\exo{4}{Moyen}{Cardinal et parties d'un ensemble}
Soit $E=\{a,b,c\}$. Énumérer tous les sous-ensembles de $E$, donner le cardinal de $\mathcal P(E)$ et comparer à $2^{|E|}$.

\exo{5}{Moyen}{Une inclusion stricte}
Soit $\mathbb N=\{0,1,2,\ldots\}$ et
$A=\{n\in\mathbb N:n\le5\}$, $B=\{1,2,3,4,5\}$.
Comparer $A$ et $B$ : établir une inclusion puis décider si les ensembles sont égaux.

\section{Opérations sur les ensembles et lois de De Morgan}

\exo{6}{Facile}{Union, intersection et différence}
Pour $A=\{1,2,3,4\}$ et $B=\{3,4,5\}$, déterminer
$A\cup B$, $A\cap B$, $A\setminus B$ et $B\setminus A$.

\exo{7}{Facile}{Complémentaire dans un univers}
Dans l'univers $E=\{\mathrm C,\mathrm H,\mathrm O,\mathrm N,\mathrm S\}$,
soit $A=\{\mathrm C,\mathrm H,\mathrm O\}$.
Calculer $A^c=E\setminus A$, puis comparer ce résultat au complémentaire de $A$ dans l'univers $E'=E\cup\{\mathrm{Cl}\}$.

\exo{8}{Moyen}{Distributivité}
Pour $A=\{1,2,4\}$, $B=\{2,3\}$ et $C=\{1,3,4\}$,
calculer séparément $A\cap(B\cup C)$ et $(A\cap B)\cup(A\cap C)$ puis comparer.

\exo{9}{Moyen}{De Morgan sur des ensembles finis}
Avec $E=\{1,2,3,4,5,6\}$, $A=\{1,2,4\}$ et $B=\{2,3,5\}$, calculer les deux membres de
$(A\cup B)^c=A^c\cap B^c$ et de
$(A\cap B)^c=A^c\cup B^c$.

\exo{10}{Avancé}{Une preuve par appartenance}
Démontrer, pour des ensembles quelconques $A,B,C$ inclus dans un univers fixé :
\[
A\cap(B\cup C)=(A\cap B)\cup(A\cap C).
\]
Utiliser un élément arbitraire $x$ et une chaîne d'équivalences.

\section{Produit cartésien et contraintes de formulation}

\exo{11}{Facile}{Former un produit cartésien}
Soient $A=\{1,2\}$ et $B=\{p,q,r\}$.
Énumérer $A\times B$ et $B\times A$ ; expliquer l'importance de l'ordre et vérifier les cardinaux.

\exo{12}{Moyen}{Produit d'un ensemble par lui-même}
Pour $A=\{0,1,2\}$, représenter $A\times A$ dans un repère. Combien de couples ont des coordonnées égales ?

\exo{13}{Application}{Composition admissible d'un mélange}
Deux fractions molaires $x,y$ sont admissibles si
$x\ge0$, $y\ge0$, $x+y\le1$.
Décrire l'ensemble $D\subset\mathbb R^2$, le représenter graphiquement et tester les couples $(0{,}3;0{,}4)$ et $(0{,}8;0{,}5)$.

\exo{14}{Moyen}{Sous-ensemble de couples}
Décrire géométriquement
$S=\{(x,y)\in[0,1]^2:x+y=1\}$.
Donner trois couples de $S$ et trois couples de $[0,1]^2$ qui n'appartiennent pas à $S$.

\exo{15}{Moyen}{Relation et application}
Soit $A=\{1,2,3\}$, $B=\{a,b\}$ et
$R=\{(1,a),(2,b),(3,b)\}\subset A\times B$.
Expliquer pourquoi $R$ définit une application $f:A\to B$ ; donner ses images et ses antécédents.

\section{Applications, injection, surjection et réciproque}

\exo{16}{Facile}{Image et antécédent}
Soit $f:\mathbb R\to\mathbb R$, $f(x)=3x-2$.
Calculer les images de $-1$, $0$ et $2$ ; déterminer l'antécédent de $7$.

\exo{17}{Moyen}{Tester l'injectivité}
Soit $g:\mathbb R\to\mathbb R$, $g(x)=x^2$.
Trouver deux antécédents différents d'une même image et conclure sur l'injectivité. Étudier ensuite la restriction $g:[0,+\infty[\to[0,+\infty[$.

\exo{18}{Moyen}{L'ensemble d'arrivée compte}
Étudier la surjectivité de $h(x)=x^2$ d'abord comme application $\mathbb R\to\mathbb R$, puis comme application $\mathbb R\to[0,+\infty[$. Comparer les résultats.

\exo{19}{Moyen}{Bijection affine}
Pour $f:\mathbb R\to\mathbb R$, $f(x)=2x+5$, démontrer l'injectivité et la surjectivité. Déterminer $f^{-1}$ et vérifier les deux compositions avec l'identité.

\exo{20}{Moyen}{L'ordre de composition}
Soit $f(x)=2x+1$ et $g(x)=x^2$ de $\mathbb R$ vers $\mathbb R$.
Calculer $g\circ f$ et $f\circ g$, puis leurs valeurs en $x=1$. En déduire pourquoi la composition n'est pas commutative en général.

\section{Logique, quantificateurs et preuves}

\exo{21}{Facile}{Connecteurs}
Construire les tables de vérité de $P\land Q$, $P\lor Q$ et $\neg P$.
Expliquer par une phrase la différence entre « et », « ou » et « non ».

\exo{22}{Moyen}{Implication et équivalence}
Pour des propositions $P,Q$, établir la table de vérité de $P\Rightarrow Q$ et de $P\Leftrightarrow Q$.
Donner un exemple montrant que l'implication n'est pas toujours une équivalence.

\exo{23}{Moyen}{Négation des quantificateurs}
Écrire les négations de :
$\forall x\in\mathbb R,\ x^2\ge0$ et
$\exists n\in\mathbb Z,\ n^2=2$.
Préciser la valeur de vérité des propositions initiales.

\exo{24}{Moyen}{Preuve directe et contraposée}
Démontrer que si un entier $n$ est pair, alors $n^2$ est pair. Énoncer et démontrer la contraposée : si $n^2$ est impair, alors $n$ est impair.

\exo{25}{Avancé}{Contre-exemple et ensemble de validité}
Étudier l'affirmation « Pour tout réel $x$, $x^2\ge x$ ».
Fournir, si nécessaire, un contre-exemple et déterminer l'ensemble des $x$ pour lesquels cette inégalité est vraie.

\section{Applications aux données et aux contraintes chimiques}

\exo{26}{Application}{Collections d'espèces chimiques}
Un premier dosage suit les espèces $A=\{\mathrm{Fe},\mathrm O,\mathrm H\}$, un second suit $B=\{\mathrm{Fe},\mathrm C,\mathrm O\}$.
Calculer $A\cup B$, $A\cap B$ et $B\setminus A$ ; expliquer leur signification expérimentale.

\exo{27}{Application}{Deux critères de qualité}
Pour un lot, $P$ signifie « la concentration est conforme » et $Q$ signifie « le pH est conforme ».
Traduire en logique « le lot est accepté seulement si les deux critères sont satisfaits », puis écrire la négation et vérifier un lot satisfaisant $P$ mais pas $Q$.

\exo{28}{Application}{Domaine d'une formulation}
Un mélange contient deux constituants aux proportions $x$ et $y$ dans
$D=\{(x,y)\in\mathbb R^2:x\ge0,\ y\ge0,\ x+y\le1\}$.
Pour $x=0{,}2$, déterminer les valeurs possibles de $y$. Représenter l'ensemble des couples admissibles et expliquer le sens du troisième constituant de proportion $1-x-y$.

\exo{29}{Application}{Une calibration comme application}
Un signal $s$ est associé à une concentration $c$ par
$f:[0,+\infty[\to[1,+\infty[$, $f(c)=2c+1$.
Montrer que $f$ est bijective ; déterminer $f^{-1}$ et interpréter le calcul de $c$ à partir d'un signal $s=9$.

\exo{30}{Synthèse}{Trois échantillons et deux détecteurs}
Soit $E=\{A,B,C\}$ l'ensemble des échantillons et $F=\{d_1,d_2\}$ l'ensemble des détecteurs.
L'application $f:E\to F$ est définie par
$f(A)=d_1$, $f(B)=d_1$, $f(C)=d_2$.
\begin{enumerate}
\item Écrire son graphe sous forme d'ensemble de couples ; déterminer $f(E)$ et chaque ensemble d'antécédents.
\item Dire si $f$ est injective, surjective, bijective ; justifier par les définitions.
\item Traduire avec des quantificateurs « chaque détecteur reçoit au moins un échantillon » et écrire sa négation.
\end{enumerate}

\section*{Auto-évaluation}
\addcontentsline{toc}{section}{Auto-évaluation}
\begin{objectifbox}
Je sais calculer des opérations ensemblistes, démontrer une identité, former un produit cartésien, préciser les images et antécédents, tester injection/surjection/bijection, composer des applications et raisonner avec quantificateurs et contre-exemples.
\end{objectifbox}
\end{document}

